% Part: turing-machines 
% Chapter: machines-computations 
% Section: introduction

\documentclass[../../../include/open-logic-section]{subfiles}

\begin{document}

\olfileid{tur}{mac}{int}
\olsection{引言}

一个函数（比如从 $\Nat$ 到 $\Nat$ 的函数）是\emph{可计算的}，这是什么意思？最早的、也是最著名的答案之一是：如果一个函数能被图灵机计算，那么它就是可计算的。这个概念由 图灵于 1936 年提出。图灵机是\emph{计算模型}的一个例子——它用数学上精确的方式定义了“计算过程”这一概念。这种定义的确切含义尚有争议，但人们普遍同意，图灵机是规定计算过程的一种方式。尽管“图灵机”这个术语会让人联想到带有活动部件的物理机器，但严格来说，图灵机是一种纯粹的数学构造，因而它是对计算过程概念的一种理想化。例如，我们对图灵机的时间或内存需求没有任何限制：即使某项计算所需的存储空间或步骤比宇宙中的原子还要多，图灵机依然能完成计算。

\begin{explain}
或许最好把图灵机看作一种特殊假想机制的程序。这个机制由一个\emph{纸带}和一个\emph{读写头}组成。在我们这个版本的图灵机中，纸带是单向（向右）无限的，并且被划分成\emph{方格}，每个方格可以包含一个来自有限\emph{字母表}的符号。这样的字母表可以包含任意多个不同的符号，但我们主要使用三个符号：$\TMendtape$、$\TMblank$ 和 $\TMstroke$。当机制启动时，除了最左边的方格包含 $\TMendtape$ 以及有限个包含\emph{输入}的方格外，纸带的其余部分是空的（即每个方格包含符号 $\TMblank$）。在任何时刻，该机制都处于有限个\emph{状态}之一。开始时，读写头扫描最左边的方格，并处于一个指定的\emph{初始状态}。在该机制运行的每一步，当前扫描方格的内容、机制所处的状态以及图灵机程序共同决定了接下来会发生什么。图灵机程序由一个偏函数给出，该函数以状态~$q$ 和符号~$\sigma$ 作为输入，并输出一个三元组~$\tuple{q', \sigma', D}$。每当机制处于状态 $q$ 并读取符号 $\sigma$ 时，它就会将当前方格上的符号替换为 $\sigma'$，读写头根据 $D$ 为 $\TMleft$、$\TMright$ 或 $\TMstay$ 而向左移动、向右移动或保持不动，并且机制进入状态~$q'$。

例如，考虑 \olref{fig:tm} 中的情形。
\begin{figure}
  \olcachedasset{turing-machine}
  \caption{一台正在执行其程序的图灵机。}
  \ollabel{fig:tm}
\end{figure}
图灵机纸带的可见部分在最左边的方格包含纸带结束符 $\TMendtape$，随后是三个 $1$、一个 $0$ 以及另外四个 $1$。读写头正在从左数第三个方格读取内容，该方格包含一个~$1$，且机器处于状态~$q_1$——我们称“机器正在状态~$q_1$ 下读取一个 $1$”。如果图灵机的程序在输入 $\tuple{q_1, 1}$ 时返回三元组 $\tuple{q_2, 0, \TMstay}$，那么机器就会将第三个方格上的~$1$ 替换为~$0$，让读写头停留在原地，并切换到状态~$q_2$。然后，如果程序在输入 $\tuple{q_2, 0}$ 时返回 $\tuple{q_3, 0, \TMright}$，那么机器就会用另一个~$0$ 覆盖原来的~$0$（实际上，保持读写头下方纸带内容不变），向右移动一个方格，并进入状态~$q_3$。以此类推。

我们说，当机器处于某个状态 $q_n$ 并读到某个符号 $\sigma$，而程序中没有针对 $\tuple{q_n, \sigma}$ 的指令时，即针对输入 $\tuple{q_n,\sigma}$ 的转移函数无定义时，机器\emph{停机}。换句话说，机器没有要执行的指令，此时它停止运行。有时停机用一个特定的停机状态~$h$ 来表示。这一点将在后面更详细地展示。
\end{explain}

\begin{digress}
图灵论文《论可计算数》的精妙之处在于，他不仅提出了一个形式定义，还论证了这个定义捕捉了直观的可计算性概念。
从定义出发，不难看出任何由图灵机可计算的函数在直观意义上也是可计算的。
图灵提供了三类论证来说明其逆亦真，即任何我们自然会视作可计算的函数，都能由这样的机器计算。用图灵的话说，它们是：
\begin{enumerate}
\item 直接诉诸直觉。
\item 证明两个定义的等价性（如果新定义具有更强的直观说服力）。
\item 给出可计算数的大类作为例子。
\end{enumerate}
我们的目标是尝试定义“原则上”的可计算性概念，即不考虑时间和空间的实际限制。
当然，有了最宽泛的可计算性定义之后，人们就可以进而研究资源受限的计算；这构成了被称为“计算复杂性”的学科的核心。
\end{digress}

\begin{history}
Alan 图灵 在 1936 年发明了图灵机。虽然他当时的兴趣是一阶逻辑的可判定性，但这篇论文被描述为关于计算机设计基础的权威文献。在文中，图灵主要关注可计算实数，即其十进制展开是可计算的实数；但他指出，将他的概念推广到自然数上的可计算函数等并不困难。请注意，这比第一台实际工作的通用计算机的诞生（1941年由德国人Konrad Zuse在他父母的客厅里建造）早了整整五年，比 Alan 图灵 及其在布莱切利园的同事们建造破译密码的巨人计算机（Colossus，1943年）早了七年，比美国的ENIAC（1945年）早了九年，比第一台英国通用计算机——曼彻斯特小型实验机（SSEM）在曼彻斯特建成（1948年）早了十二年，比美国人首次测试BINAC（1949年）早了十三年。曼彻斯特SSEM的独特之处在于它是第一台存储程序计算机——此前的机器必须为每个新任务手动重新接线。
\end{history}

\end{document}
